% English translation of checked_pages/esq_119_045.tex; source: cote 119, site 45, PDF page 46, typescript page 43.
% AI-translated by GPT-6 Astra (OpenAI, ChatGPT Pro mode).
\scanpage{119/45}
\mspage{43}

volumes to be published by Hermann, under the common title ``\uline{Pursuing}
\uline{Stacks}''. I currently expect (after successive expansions
of the initial undertaking) that the manuscript\frapparatus{V-FR-045-001} of the two
volumes together, which I hope to finish during the year so as not
to have to return to it, will run to about 1,500 typewritten pages.
These two volumes are moreover, for me, the first in a
larger series, under the common title ``\uline{Mathematical Reflections}'',
in which I intend to develop to some extent certain of the themes sketched
in the present report.

Since this is a work currently being written,
and indeed approaching completion, whose first volume will probably appear
this year and will contain a detailed introduction, it
is probably less interesting for me to dwell here on this theme
of reflection, and I shall therefore be content to speak of it very briefly.
This work seems to me somewhat marginal in relation to the themes
I have just sketched, and does not represent (it seems to me)
a real renewal of outlook or approach in relation to
my mathematical interests and vision before 1970. If I
suddenly resolved to undertake it, this was almost out of desperation,
when nearly twenty years have elapsed since a number of visibly
fundamental questions, ripe to be carried through to completion,
were posed in quite clear terms, without
anyone seeing them or taking the trouble to probe them. Even
today, the basic structures involved in the
homotopy viewpoint in topology, including even commutative
homological algebra, are not understood, and to my
knowledge, after the work of Verdier, Giraud, and Illusie on
this theme (which constitute so many ``kick-offs'' still
awaiting a continuation ...), there has been no effort in this direction. I
should probably make an exception for\unsure{.}\tnote{TN-045-01}{A visible baseline dot follows ``pour'' in the scan. Its punctuation status remains uncertain; the mark is retained after ``for'' here.} the axiomatization
carried out by Quillen on the notion of a model category, in
the late 60s, and taken up in various forms by
various authors. This work, at the time and still today, has
greatly attracted me and taught me much, while going in a rather

\margin{45}
